% Owner: Samriddhi (Member 4)
\subsection{Verification approach}
The RTL is Verilog-2001 \cite{ieee1364}, simulated with Icarus Verilog \cite{icarus}. Every
testbench checks its own results and prints \texttt{PASS} or \texttt{FAIL}, and one
\sig{Makefile} runs the whole flow (\texttt{make all}).

\textbf{Assembler and programs.} The Phase~I two-pass assembler now supports the new
instructions and the pseudo-instructions \texttt{j}, \texttt{ret} and \texttt{mv}. Its
self-test checks 29 encodings against hand-computed values, including well-known reference
encodings such as \texttt{ret}\,=\,\texttt{0x00008067}. Four new programs, each with different
memory behaviour, join the Phase~I four:
\begin{itemize}
  \item \sig{bubble\_sort}: sorts 16 signed words. It has strong temporal locality and issues
        stores in pairs, one pair per swap.
  \item \sig{subroutine}: multiply, GCD and popcount subroutines and a nested call through a
        stack. Its 59-word code footprint exceeds an 8-line instruction cache.
  \item \sig{transpose}: an $8\times8$ matrix transpose. Its reads are sequential (spatial
        locality) and its writes have a stride of 8 words.
  \item \sig{isa\_test}: every new instruction once, with hand-computed results.
\end{itemize}
The final register and memory values of every program were worked out by hand (and with plain
Python) and are checked on both Phase~I systems.

\textbf{Golden model.} \sig{tools/iss.py} is an instruction-set simulator written independently
of the RTL and the assembler. It decodes the instruction bits itself. It predicts the final
registers, data memory and \sig{out\_port}; the instruction, load and store counts; and the
misses of split and unified direct-mapped caches of any size, with write-through and
no-write-allocate.

\textbf{Cache unit test.} \sig{tb\_cache} drives four caches with memories of latency 3 or 1.
Its 41 checks cover miss timing ($1+L$ cycles), hits, a buffered store, a store waiting for a
busy buffer, no-write-allocate, a read after a buffered store miss, conflict eviction, store
hits with and without a buffer, and uncached mode. We confirmed that the test can fail by
running it against three deliberately broken copies of the RTL: one with the read-after-write
wait removed, one where a store miss writes the line, and one with an off-by-one memory latency.
Each copy fails at least one check (\texttt{make mutate}).

\textbf{System testbench.} \sig{tb\_system} runs one program on five systems at once: the ideal
Harvard reference, the modified Harvard system, the unified-cache system, the no-cache system
and the Phase~I baseline. It checks that:
\begin{itemize}
  \item registers, all 256 data words, \sig{out\_port}, instruction counts, data-access counts
        and fault cycles match the reference;
  \item the identities (\ref{eq:mh})--(\ref{eq:bus}) hold;
  \item there are no conflicts on a single bus, and no conflicts or write-buffer stalls when the
        buffer is off;
  \item with $L=1$, the no-cache system equals the Phase~I baseline in cycle count and in bus
        address and FSM state in every cycle.
\end{itemize}

\textbf{Sweep.} \sig{tools/sweep.py} runs \sig{tb\_system} for eight programs
$\times$ $L\in\{1,2,4,8,16\}$ $\times$ $N\in\{1,2,4,8,16,32\}$ lines per split cache
$\times$ write buffer off/on, plus a placement experiment (Section~\ref{sec:split}): 528 runs in
all. It cross-checks every RTL miss count and instruction count against the golden model, and
every final state against the ISS. With the buffer off it also checks an exact stall model.
Cached systems must satisfy $S_I = L\cdot M_I$ and $S_D = L\cdot M_D + (L-1)\cdot W$. The
no-cache system must satisfy $S_I = (L-1)\cdot I$ and $S_D=(L-1)(R+W)$. Here $M$ is a miss count
and $R$, $W$ are the loads and stores to RAM. All 528 runs pass, with 9888 golden-model checks
and no mismatch. The script also writes the CSV file, the summary tables and the plot data used
below, so no number in this report was typed by hand.
